\documentclass[10pt, a4paper]{report}
\usepackage[polish]{babel}
\usepackage[utf8]{inputenc}
\usepackage{graphicx} % Required for inserting images
\usepackage{titlesec}
\usepackage{tocloft}
% setting up the mirror margins
\usepackage[left=2.5cm,right=2.5cm,top=2.5cm,bottom=2.5cm,bindingoffset=1cm]{geometry}

%SETTING FONT TO ARIAL
\usepackage{helvet}
\renewcommand{\familydefault}{\sfdefault}

% for checking the font sizes to make sure they match what the template wants
\makeatletter
\newcommand\thefontsize{The current font size is: \f@size pt}
\makeatother

% TITLE PAGE IS A SEPERATE FILE GENERATED FROM THE MOJAPG SYSTEM

% setting up paragraph indent size
\setlength\parindent{1cm}

\begin{document}

% setting sizes and format for sections-subsubsection
\setcounter{secnumdepth}{3}
\setcounter{tocdepth}{3}
\renewcommand{\thesection}{\arabic{section}.}
\renewcommand{\thesubsection}{\arabic{section}.\arabic{subsection}.}
\renewcommand{\thesubsubsection}{\arabic{section}.\arabic{subsection}.\arabic{subsubsection}.}


\titleformat{\section} {\fontsize{12pt}{15pt}\selectfont\bfseries} {\thesection}{1em}{}
\titleformat{\subsection} {\fontsize{10pt}{13pt}\selectfont\bfseries\itshape} {\thesubsection}{1em}{}
\titleformat{\subsubsection} {\fontsize{10pt}{13pt}\selectfont\itshape} {\thesubsubsection}{1em}{}

% ---------------------------------------------------------------------------------
% STRESZCZENIE / ABSTRACT
% ---------------------------------------------------------------------------------
\section*{STRESZCZENIE} % needs to be 12pts, odstep przed 12pt, odstep po 6pt

% W Streszczeniu nie piszemy co znajduje się w konkretnych rozdziałach.
% Rolą Streszczenia jest ogólne nakreślenie szerszej tematyki i problematyki podjętej w pracy dyplomowej oraz tego co w pracy wykonano.
% Streszczenie powinno zawierać określenie problemu naukowego lub praktycznego do rozwiązania, cel i zakres pracy, zastosowane metody badań, wyniki i najważniejsze wnioski
% Warto, aby w Streszczeniu zawrzeć krótki zarys tematu (kilka zdań).

Projekt inżynierski...


\noindent
\textbf{Słowa kluczowe}:

\noindent
\textbf{Dziedzina nauki i techniki, zgodnie z wymogami OECD}:\\
Nauki o komputerach i informatyka

\newpage
\section*{ABSTRACT} % same as what is in strreszczenie just on a seperate page and in english

\noindent
\textbf{Keywords}:

\noindent
\textbf{Field of science and technology in accordance with OECD requirements}:\\
Computer Science and Information technology
% ---------------------------------------------------------------------------------
% SPIS TREŚCI / TABLE OF CONTENTS
% ---------------------------------------------------------------------------------
\newpage
% TODO: SET INTERLINIA
% TODO: SET ODPTEP MIEDZY AKAPITAMI

% W spisie treści stosujemy pojedynczą interlinię
% odstęp akapitu przed równy 0 pkt., natomiast odstęp po równy 6 pkt.
% Wysunięcie akapitów z tytułami rozdziałów może wynosić 0,75 cm, a dla podrozdziałów 1,5 cm
\addcontentsline{toc}{section}{SPIS TREŚCI}
\renewcommand{\contentsname}{\fontsize{12pt}{15pt}\selectfont\bfseries SPIS TREŚCI}

 % TODO: THESE INDENTS ARE COMPLETELY EYEBALLED. IDK HOW TO MAKE THEM BETTER NGL
\cftsetindents{section}{0em}{1.75em} % no indentation
\cftsetindents{subsection}{1.5em}{2.5em} % no indentation
\cftsetindents{subsubsection}{4em}{3.5em} % no indentation

\tableofcontents

% ---------------------------------------------------------------------------------
% WYKAZ SKROTOW / TABLE OF MOST IMPORTANT KEYWORDS ETC
% ---------------------------------------------------------------------------------

\newpage
\addcontentsline{toc}{section}{Wykaz ważniejszych oznaczeń i skrótów}
\section*{Wykaz ważniejszych oznaczeń i skrótów}


% pojedynczą interlinię, odstęp przed równy 0 pkt., odstęp po równy 6 pkt. Symbol wyrównać do lewej, następnie dać tabulację, myślnik, znak tabulacji i opis symbolu.
% Opisu nie kończymy kropką, przecinkiem ani średnikiem.

% ---------------------------------------------------------------------------------
% WŁAŚCIWA CZĘŚĆ PRACY / PAPER ITSELF
% ---------------------------------------------------------------------------------

\section{teeeeeeeeeesses}
\subsection{teeeeeeeeest}
\subsubsection{}

\section{}
\section{}
\section{}
\section{}
\section{}
\section{}
\section{}
\section{Podsumowanie}


% ---------------------------------------------------------------------------------
% WYKAZ ŹRÓDEŁ / BIBLIOGRAPHY ETC
% ---------------------------------------------------------------------------------

\newpage
\addcontentsline{toc}{section}{WYKAZ LITERATURY}
\section*{WYKAZ LITERATURY}
% stosujemy pojedynczą interlinię oraz odstęp przed równy 0 pkt, natomiast odstęp po równy 6 pkt.
% Można zastosować wysunięcie równe 1 cm.
% wykonujemy w porządku alfabetycznym względem nazwiska pierwszego autora/twórcy.
% W przypadku wielu źródeł, których autorem jest jeden twórca, ustawiamy je w kolejności od najstarszego do najnowszego.

% setting the list numbers as [x] instrad of regular x.x
\renewcommand{\labelenumi}{[\arabic{enumi}]}

% TODO: SET THE PROPER INDENTS FOR THIS TOO!
\begin{enumerate}
    \item ZRODLO 1 ETC
\end{enumerate}

\newpage
\addcontentsline{toc}{section}{WYKAZ RYSUNKÓW}
\section*{WYKAZ RYSUNKÓW}

% Kolejność w spisie odpowiada kolejności numeracji rysunków, tabel, załączników.
% Stosujemy pojedynczą interlinię oraz odstęp przed równy 0 pkt, natomiast odstęp po równy 6 pkt.

% Akapit rozpoczyna się numerem rysunku/tabeli, następnie umieszczamy tabulator, tytuł rysunku/tabeli, tabulator z opcją podkreślenia punktami z wyrównaniem do prawej, na końcu numer strony, na której znajduje się rysunek/tabela.

\newpage
\addcontentsline{toc}{section}{WYKAZ TABEL}
\section*{WYKAZ TABEL}

\newpage
\addcontentsline{toc}{section}{WYKAZ LISTINGÓW}
\section*{WYKAZ LISTINGÓW}


\end{document}
